\chapter{Graph Traversal}
\chapterepigraph{Every graph problem starts with a traversal. BFS and DFS
are the same idea -- visit a node, then visit its unvisited neighbours --
differing only in whether a queue or a stack (recursion) decides who is
next. BFS's level-by-level order gives shortest paths in unweighted graphs;
DFS's dive-deep order exposes components, cycles, and structure. This
chapter is the foundation the rest of the book stands on.}

\section{BFS and DFS: One Idea, Two Orders}
\label{sec:cses-graph-traversal-intro}

A traversal marks nodes visited so none is processed twice, and explores
outward from a source. The only choice is the order neighbours are taken.

\begin{ideabox}[Queue for Shortest, Stack for Structure]
\textbf{BFS} uses a queue: it visits all nodes at distance $1$, then $2$,
and so on, so the first time it reaches a node is via a shortest
(fewest-edges) path. \textbf{DFS} uses recursion (an implicit stack): it
follows one path as deep as possible before backtracking, which naturally
reveals connected components, cycles (a back edge to an ancestor), and
bipartiteness (a 2-colouring conflict).
\end{ideabox}

\begin{patternbox}[What Each Problem Here Needs]
\begin{itemize}
  \item \textbf{Count/flood connected regions} -- DFS or BFS flood fill:
        Counting Rooms.
  \item \textbf{Shortest path / reconstruct path in an unweighted graph
        or grid} -- BFS: Labyrinth, Message Route.
  \item \textbf{Is everything connected / how many edges to connect} --
        component counting: Building Roads.
  \item \textbf{2-colour the graph} -- BFS/DFS bipartite check: Building
        Teams.
  \item \textbf{Find a cycle} -- DFS back edge (or DSU) in an undirected
        graph: Round Trip.
  \item \textbf{Multi-source spread / escape} -- BFS from many sources at
        once: Monsters.
\end{itemize}
\end{patternbox}

Each section gives the CSES-style input handling, the traversal, path or
answer reconstruction where required, a dry run, complete C++, and
complexity.

\newpage

\section{Counting Rooms}
\label{sec:cses-counting-rooms}

\problemstatement{Given an $n\times m$ grid of floor (\texttt{.}) and wall
(\texttt{\#}) cells, count the number of rooms: maximal groups of floor
cells connected through shared edges (four directions).}

\begin{patternbox}
\emph{``Count connected regions in a grid''} is a \textbf{flood fill}:
whenever an unvisited floor cell is found, start a traversal that marks its
entire room, and increment a counter. Each new traversal is exactly one
room.
\end{patternbox}

\subsection*{Flood fill counts components}

Scan every cell. When you hit a floor cell not yet visited, it belongs to a
new room, so increment the room count and run a BFS or DFS that marks all
floor cells reachable from it (moving only up/down/left/right through
floor). Because each cell is marked the first time it is reached and never
revisited, the number of traversals launched equals the number of rooms.

\begin{ideabox}[One Launch per Component]
The visited grid guarantees each cell is consumed once. So the count of
``fresh starts'' -- unvisited floor cells that trigger a new flood -- is the
component count. This is the grid form of connected-component counting.
\end{ideabox}

\subsection*{Algorithm}

\begin{enumerate}
  \item For each cell $(i,j)$: if it is floor and unvisited, increment the
        room count and flood-fill from it.
  \item The flood marks every floor cell reachable through the four
        directions.
  \item Output the room count.
\end{enumerate}

\subsection*{Dry run}

\dryrun{$3\times3$ grid: row0 ``.\#.'', row1 ``.\#.'', row2 ``...''.}

\begin{itemize}
  \item Start at $(0,0)$: flood the left column and bottom row -- one
        room.
  \item $(0,2)$ connects down to $(1,2)$ then to the bottom row, already
        visited -- so it is the \emph{same} room via the bottom.
  \item Here everything floor is one connected room, so the answer is
        $1$.
\end{itemize}

\subsection*{C++ implementation}

\begin{lstlisting}
#include <bits/stdc++.h>
using namespace std;

int n, m;
vector<string> g;
vector<vector<bool>> vis;

void bfs(int si, int sj) {
    queue<pair<int,int>> q;
    q.push({si, sj}); vis[si][sj] = true;
    int dr[] = {-1,1,0,0}, dc[] = {0,0,-1,1};
    while (!q.empty()) {
        auto [r, c] = q.front(); q.pop();
        for (int k = 0; k < 4; k++) {
            int nr = r + dr[k], nc = c + dc[k];
            if (nr < 0 || nr >= n || nc < 0 || nc >= m) continue;
            if (g[nr][nc] == '#' || vis[nr][nc]) continue;
            vis[nr][nc] = true;
            q.push({nr, nc});
        }
    }
}

int main() {
    cin >> n >> m;
    g.resize(n);
    for (auto& row : g) cin >> row;
    vis.assign(n, vector<bool>(m, false));

    int rooms = 0;
    for (int i = 0; i < n; i++)
        for (int j = 0; j < m; j++)
            if (g[i][j] == '.' && !vis[i][j]) { rooms++; bfs(i, j); }

    cout << rooms << "\n";
    return 0;
}
\end{lstlisting}

\begin{complexitybox}
\textbf{Time Complexity.} $O(nm)$ -- each cell is visited once.
\textbf{Space Complexity.} $O(nm)$ for the visited grid and BFS queue.
\end{complexitybox}

\begin{takeawaybox}
Counting rooms is connected-component counting on an implicit grid graph:
one flood fill per unvisited region. This is the simplest graph traversal
and the template for every grid problem that follows.
\end{takeawaybox}

\section{Labyrinth}
\label{sec:cses-labyrinth}

\problemstatement{Given a grid with a start \texttt{A}, a target
\texttt{B}, floor \texttt{.}, and walls \texttt{\#}, find a shortest path
from \texttt{A} to \texttt{B} moving in the four directions. Output
\texttt{YES}, its length, and the sequence of moves (\texttt{U/D/L/R}), or
\texttt{NO}.}

\begin{patternbox}
A grid where each move costs one step is an unweighted graph, so
\textbf{BFS} finds the shortest path, and remembering \emph{how} each cell
was first reached (which move and from where) reconstructs the actual path.
\end{patternbox}

\subsection*{BFS plus per-cell parent}

Run BFS from \texttt{A}. The first time BFS reaches a cell, it does so by a
shortest path, so record the move (\texttt{U/D/L/R}) used to enter it and
its predecessor. When \texttt{B} is reached, walk the parent pointers back
to \texttt{A}, collecting the moves, then reverse them. If BFS finishes
without reaching \texttt{B}, output \texttt{NO}.

\begin{ideabox}[Store the Move That First Reached Each Cell]
Distances alone say how long the path is; storing the direction taken into
each cell lets you replay the path. Because BFS's first arrival is optimal,
these stored moves spell out one shortest route.
\end{ideabox}

\subsection*{Algorithm}

\begin{enumerate}
  \item BFS from \texttt{A}; for each newly visited cell store the move
        character that entered it.
  \item On reaching \texttt{B}, backtrack via parents, collect moves,
        reverse.
  \item Output \texttt{YES}, the length, and the move string; else
        \texttt{NO}.
\end{enumerate}

\subsection*{Dry run}

\dryrun{\texttt{A} at $(0,0)$, \texttt{B} at $(0,2)$, top row open.}

\begin{itemize}
  \item BFS reaches $(0,1)$ via \texttt{R}, then $(0,2)$ via \texttt{R}.
  \item Backtrack from \texttt{B}: moves \texttt{R},\texttt{R}; reversed
        gives \texttt{RR}, length $2$.
\end{itemize}

\subsection*{C++ implementation}

\begin{lstlisting}
#include <bits/stdc++.h>
using namespace std;

int main() {
    int n, m;
    cin >> n >> m;
    vector<string> g(n);
    for (auto& row : g) cin >> row;

    int sr, sc, er, ec;
    for (int i = 0; i < n; i++)
        for (int j = 0; j < m; j++) {
            if (g[i][j] == 'A') { sr = i; sc = j; }
            if (g[i][j] == 'B') { er = i; ec = j; }
        }

    vector<vector<char>> move(n, vector<char>(m, 0));
    vector<vector<bool>> vis(n, vector<bool>(m, false));
    int dr[] = {-1,1,0,0}, dc[] = {0,0,-1,1};
    char mv[] = {'U','D','L','R'};

    queue<pair<int,int>> q;
    q.push({sr, sc}); vis[sr][sc] = true;
    while (!q.empty()) {
        auto [r, c] = q.front(); q.pop();
        for (int k = 0; k < 4; k++) {
            int nr = r + dr[k], nc = c + dc[k];
            if (nr < 0 || nr >= n || nc < 0 || nc >= m) continue;
            if (g[nr][nc] == '#' || vis[nr][nc]) continue;
            vis[nr][nc] = true;
            move[nr][nc] = mv[k];
            q.push({nr, nc});
        }
    }

    if (!vis[er][ec]) { cout << "NO\n"; return 0; }
    string path;
    int r = er, c = ec;
    while (!(r == sr && c == sc)) {
        char mvc = move[r][c];
        path += mvc;
        if (mvc == 'U') r++; else if (mvc == 'D') r--;
        else if (mvc == 'L') c++; else c--;      // step back
    }
    reverse(path.begin(), path.end());
    cout << "YES\n" << path.size() << "\n" << path << "\n";
    return 0;
}
\end{lstlisting}

\begin{complexitybox}
\textbf{Time Complexity.} $O(nm)$ -- one BFS over the grid.
\textbf{Space Complexity.} $O(nm)$ for visited and move grids.
\end{complexitybox}

\begin{takeawaybox}
Shortest path on an unweighted grid is BFS; storing the entering move per
cell turns the distance into a printable route. The backtrack-and-reverse
pattern reappears in every ``print the path'' problem, including Message
Route next.
\end{takeawaybox}

\section{Message Route}
\label{sec:cses-message-route}

\problemstatement{Given an undirected graph of $n$ computers and $m$
connections, find a route with the fewest edges from computer $1$ to
computer $n$. Output the number of computers on the route and the route
itself, or \texttt{IMPOSSIBLE}.}

\begin{patternbox}
Fewest edges in an unweighted graph is \textbf{BFS}. Recording each node's
predecessor lets the shortest route be reconstructed by walking parents
back from $n$ to $1$.
\end{patternbox}

\subsection*{BFS with parent pointers}

BFS from node $1$. The first time a node is dequeued-into, it is via a
shortest route, so store the node it came from. When $n$ is reached, follow
$\textit{parent}$ from $n$ back to $1$, then reverse to get the route in
forward order. If $n$ is never visited, the answer is \texttt{IMPOSSIBLE}.

\begin{ideabox}[Parents Turn Distances into Routes]
This is Labyrinth on a general graph rather than a grid: the same BFS,
the same parent-backtrack. On any unweighted graph, ``fewest edges''
plus ``print the path'' is always BFS plus predecessors.
\end{ideabox}

\subsection*{Algorithm}

\begin{enumerate}
  \item Build adjacency lists. BFS from $1$, storing $\textit{parent}[v]$
        on first visit.
  \item If $n$ unvisited, print \texttt{IMPOSSIBLE}.
  \item Else backtrack $n\to\ldots\to1$, reverse, and print the count and
        route.
\end{enumerate}

\subsection*{Dry run}

\dryrun{Edges $1\!-\!2$, $2\!-\!3$, $1\!-\!3$; $n=3$.}

\begin{itemize}
  \item BFS from $1$ reaches $3$ directly (edge $1\!-\!3$), parent
        $[3]=1$.
  \item Route $1\to3$: two computers. (The two-edge path $1\!-\!2\!-\!3$
        is longer and not chosen.)
\end{itemize}

\subsection*{C++ implementation}

\begin{lstlisting}
#include <bits/stdc++.h>
using namespace std;

int main() {
    int n, m;
    cin >> n >> m;
    vector<vector<int>> adj(n + 1);
    for (int i = 0; i < m; i++) {
        int a, b; cin >> a >> b;
        adj[a].push_back(b);
        adj[b].push_back(a);
    }

    vector<int> parent(n + 1, 0);
    vector<bool> vis(n + 1, false);
    queue<int> q;
    q.push(1); vis[1] = true;
    while (!q.empty()) {
        int u = q.front(); q.pop();
        for (int v : adj[u])
            if (!vis[v]) { vis[v] = true; parent[v] = u; q.push(v); }
    }

    if (!vis[n]) { cout << "IMPOSSIBLE\n"; return 0; }
    vector<int> route;
    for (int v = n; v != 0; v = parent[v]) route.push_back(v);
    reverse(route.begin(), route.end());

    cout << route.size() << "\n";
    for (int v : route) cout << v << " ";
    cout << "\n";
    return 0;
}
\end{lstlisting}

\begin{complexitybox}
\textbf{Time Complexity.} $O(n+m)$ -- one BFS. \textbf{Space Complexity.}
$O(n+m)$ for adjacency lists, parents, and the queue.
\end{complexitybox}

\begin{takeawaybox}
Message Route is the graph (not grid) version of shortest-path-with-route:
BFS for fewest edges, parent pointers for the path. Master this pair --
BFS distance and parent reconstruction -- as the reflex for all unweighted
shortest-path queries.
\end{takeawaybox}

\section{Building Roads}
\label{sec:cses-building-roads}

\problemstatement{Given $n$ cities and $m$ existing roads, find the minimum
number of new roads needed so that every city is reachable from every
other, and output which roads to build.}

\begin{patternbox}
\emph{``Connect all components with the fewest new edges''} reduces to
counting connected components: if there are $c$ components, exactly $c-1$
new roads suffice -- one linking a representative of each component to a
common one.
\end{patternbox}

\subsection*{Components, then bridge them}

Find the connected components (by DFS/BFS or Disjoint Set). Pick one
representative city per component. Building a road from the first
component's representative to each other representative connects everything
with $c-1$ roads, which is optimal -- no fewer can connect $c$ separate
groups. Collect and print those representative pairs.

\begin{ideabox}[c Components Need Exactly $c-1$ Links]
Each new road can reduce the number of components by at most one, so going
from $c$ components to $1$ requires at least $c-1$ roads; connecting one
representative per component to a fixed anchor achieves that bound.
\end{ideabox}

\subsection*{Algorithm}

\begin{enumerate}
  \item Find components; record one representative node per component.
  \item Output $(\text{count}-1)$ as the number of roads.
  \item For each representative after the first, print a road from the
        first representative to it.
\end{enumerate}

\subsection*{Dry run}

\dryrun{$n=4$, roads $1\!-\!2$; cities $3$ and $4$ isolated.}

\begin{itemize}
  \item Components: $\{1,2\}$, $\{3\}$, $\{4\}$ -- reps $1,3,4$.
  \item Build $1\!-\!3$ and $1\!-\!4$: two roads. Answer $2$.
\end{itemize}

\subsection*{C++ implementation}

\begin{lstlisting}
#include <bits/stdc++.h>
using namespace std;

int main() {
    int n, m;
    cin >> n >> m;
    vector<vector<int>> adj(n + 1);
    for (int i = 0; i < m; i++) {
        int a, b; cin >> a >> b;
        adj[a].push_back(b);
        adj[b].push_back(a);
    }

    vector<bool> vis(n + 1, false);
    vector<int> reps;
    for (int s = 1; s <= n; s++) {
        if (vis[s]) continue;
        reps.push_back(s);                       // one representative per component
        queue<int> q; q.push(s); vis[s] = true;
        while (!q.empty()) {
            int u = q.front(); q.pop();
            for (int v : adj[u]) if (!vis[v]) { vis[v] = true; q.push(v); }
        }
    }

    cout << reps.size() - 1 << "\n";
    for (size_t i = 1; i < reps.size(); i++)
        cout << reps[0] << " " << reps[i] << "\n";
    return 0;
}
\end{lstlisting}

\begin{complexitybox}
\textbf{Time Complexity.} $O(n+m)$ -- one traversal to find components.
\textbf{Space Complexity.} $O(n+m)$.
\end{complexitybox}

\begin{takeawaybox}
Minimum edges to connect a graph is (components $-1$); connect one
representative per component to a common anchor. Recognising ``make it all
connected'' as component counting -- solvable by traversal or Disjoint Set
-- is the reusable idea.
\end{takeawaybox}

\section{Building Teams}
\label{sec:cses-building-teams}

\problemstatement{Given $n$ players and $m$ friendships, assign each player
to one of two teams so that no two friends are on the same team. Output a
valid assignment, or \texttt{IMPOSSIBLE}.}

\begin{patternbox}
\emph{``Split into two groups so every edge crosses''} is
\textbf{2-colouring}, i.e.\ testing bipartiteness. BFS or DFS colours each
node opposite its neighbour; a conflict (an edge between same-coloured
nodes) means an odd cycle exists and the answer is \texttt{IMPOSSIBLE}.
\end{patternbox}

\subsection*{Two-colour via traversal}

Colour a start node $1$, its neighbours $2$, their neighbours $1$, and so
on, alternating on every edge. If the traversal ever tries to colour a node
the same as an already-coloured neighbour, the graph contains an odd cycle
and cannot be 2-coloured. Repeat over all components, since the graph may
be disconnected.

\begin{ideabox}[Bipartite iff No Odd Cycle]
Alternating colours along edges is consistent exactly when every cycle has
even length. The first same-colour edge encountered is a certificate of an
odd cycle, so the traversal both constructs the teams and detects
impossibility.
\end{ideabox}

\subsection*{Algorithm}

\begin{enumerate}
  \item For each uncoloured node, BFS/DFS colouring neighbours with the
        opposite colour.
  \item If a neighbour already has the same colour, output
        \texttt{IMPOSSIBLE}.
  \item Otherwise output every node's colour ($1$ or $2$).
\end{enumerate}

\subsection*{Dry run}

\dryrun{Friendships $1\!-\!2$, $2\!-\!3$; $n=3$.}

\begin{itemize}
  \item Colour $1{=}1$, $2{=}2$, $3{=}1$. No same-colour edge.
  \item Teams: player $1$ and $3$ on team $1$, player $2$ on team $2$.
\end{itemize}

\subsection*{C++ implementation}

\begin{lstlisting}
#include <bits/stdc++.h>
using namespace std;

int main() {
    int n, m;
    cin >> n >> m;
    vector<vector<int>> adj(n + 1);
    for (int i = 0; i < m; i++) {
        int a, b; cin >> a >> b;
        adj[a].push_back(b);
        adj[b].push_back(a);
    }

    vector<int> color(n + 1, 0);
    for (int s = 1; s <= n; s++) {
        if (color[s]) continue;
        color[s] = 1;
        queue<int> q; q.push(s);
        while (!q.empty()) {
            int u = q.front(); q.pop();
            for (int v : adj[u]) {
                if (!color[v]) { color[v] = 3 - color[u]; q.push(v); }
                else if (color[v] == color[u]) { cout << "IMPOSSIBLE\n"; return 0; }
            }
        }
    }
    for (int i = 1; i <= n; i++) cout << color[i] << " ";
    cout << "\n";
    return 0;
}
\end{lstlisting}

\begin{complexitybox}
\textbf{Time Complexity.} $O(n+m)$ -- one traversal. \textbf{Space
Complexity.} $O(n+m)$.
\end{complexitybox}

\begin{takeawaybox}
Two-team assignment is bipartite 2-colouring: alternate colours along
edges, and a same-colour edge proves an odd cycle. The trick $3-\text{color}$
flips between $1$ and $2$ cleanly. This underlies conflict-graph and
scheduling feasibility checks.
\end{takeawaybox}

\section{Round Trip}
\label{sec:cses-round-trip}

\problemstatement{Given an undirected graph of $n$ nodes and $m$ edges,
find any cycle (a route that starts and ends at the same node using
distinct edges) and print the sequence of nodes, or \texttt{IMPOSSIBLE}.}

\begin{patternbox}
\emph{``Find a cycle in an undirected graph''} is a \textbf{DFS} that
watches for a \textbf{back edge}: an edge to an already-visited node that
is not the immediate parent. Such an edge closes a cycle, which is
reconstructed from the DFS parent chain.
\end{patternbox}

\subsection*{Back edge closes a cycle}

Run DFS tracking each node's parent. If DFS at node $u$ finds a neighbour
$v$ that is already visited and $v$ is not $u$'s parent, then $v$ is an
ancestor of $u$ and the tree path from $v$ down to $u$, plus the edge
$u\to v$, forms a cycle. Reconstruct it by following parents from $u$ up to
$v$.

\begin{ideabox}[Not-the-Parent Is the Whole Condition]
In an undirected graph every edge appears from both ends, so seeing the
parent again is not a cycle. A visited neighbour that is \emph{not} the
parent is a genuine back edge -- the signal of a cycle. Recording parents
lets you print it.
\end{ideabox}

\subsection*{Algorithm}

\begin{enumerate}
  \item DFS from each unvisited node, storing parents.
  \item On finding a visited non-parent neighbour $v$ from $u$: walk
        parents $u\to\cdots\to v$, append $v$ again to close the loop, and
        print.
  \item If no back edge is ever found, print \texttt{IMPOSSIBLE}.
\end{enumerate}

\subsection*{Dry run}

\dryrun{Edges $1\!-\!2$, $2\!-\!3$, $3\!-\!1$.}

\begin{itemize}
  \item DFS $1\to2\to3$; from $3$, neighbour $1$ is visited and not
        $3$'s parent ($2$) -- back edge.
  \item Cycle: $1\to2\to3\to1$.
\end{itemize}

\subsection*{C++ implementation}

\begin{lstlisting}
#include <bits/stdc++.h>
using namespace std;

int n, m;
vector<vector<int>> adj;
vector<int> parent;
vector<bool> vis;
int cycleStart = -1, cycleEnd = -1;

bool dfs(int u, int p) {
    vis[u] = true;
    parent[u] = p;
    for (int v : adj[u]) {
        if (v == p) continue;                    // skip the edge back to parent
        if (vis[v]) { cycleStart = v; cycleEnd = u; return true; }
        if (dfs(v, u)) return true;
    }
    return false;
}

int main() {
    cin >> n >> m;
    adj.assign(n + 1, {});
    for (int i = 0; i < m; i++) {
        int a, b; cin >> a >> b;
        adj[a].push_back(b);
        adj[b].push_back(a);
    }
    parent.assign(n + 1, 0);
    vis.assign(n + 1, false);

    for (int s = 1; s <= n; s++)
        if (!vis[s] && dfs(s, -1)) break;

    if (cycleStart == -1) { cout << "IMPOSSIBLE\n"; return 0; }
    vector<int> cyc;
    for (int v = cycleEnd; v != cycleStart; v = parent[v]) cyc.push_back(v);
    cyc.push_back(cycleStart);
    cyc.push_back(cycleEnd);                      // close the loop
    reverse(cyc.begin(), cyc.end());

    cout << cyc.size() << "\n";
    for (int v : cyc) cout << v << " ";
    cout << "\n";
    return 0;
}
\end{lstlisting}

\begin{complexitybox}
\textbf{Time Complexity.} $O(n+m)$ -- one DFS. \textbf{Space Complexity.}
$O(n+m)$, plus recursion depth $O(n)$.
\end{complexitybox}

\begin{takeawaybox}
Undirected cycle detection is a DFS back edge to a non-parent ancestor;
parents reconstruct the cycle. (Disjoint Set also detects a cycle -- the
first edge joining two already-connected nodes -- but does not print it as
directly.)
\end{takeawaybox}

\section{Monsters}
\label{sec:cses-monsters}

\problemstatement{On a grid with your start \texttt{A}, monsters
\texttt{M}, walls \texttt{\#}, and floor \texttt{.}, you and all monsters
move one step per turn (four directions). Escape by reaching any boundary
cell before a monster occupies it. Output \texttt{YES} and a path, or
\texttt{NO}.}

\begin{patternbox}
This is a race modelled with \textbf{multi-source BFS}. First compute, for
every cell, the earliest time a monster can arrive (BFS from all monsters
at once). Then BFS from \texttt{A}, only allowed to enter a cell
\emph{strictly before} any monster could -- reaching the boundary means
escape.
\end{patternbox}

\subsection*{Two BFS passes: monsters first, then you}

Run a BFS seeded with \emph{all} monster cells simultaneously; this labels
each cell with the minimum monster arrival time. Then BFS from \texttt{A}
with your own time counter, stepping into a neighbour only if your arrival
time there is strictly less than the monster time (so a monster cannot be
there yet). Track parents to print the path. If your BFS reaches any
boundary cell, you escape; otherwise output \texttt{NO}.

\begin{ideabox}[Multi-Source BFS Gives Simultaneous Spread]
Seeding a BFS queue with every monster at once yields, in one pass, the
soonest a monster reaches each cell -- as if all monsters spread in
parallel. Comparing your arrival time against that field decides which
cells are safe to enter.
\end{ideabox}

\subsection*{Algorithm}

\begin{enumerate}
  \item Multi-source BFS from all \texttt{M} cells to get
        $\textit{monsterTime}[\,]$.
  \item BFS from \texttt{A} tracking $\textit{myTime}$ and parents;
        enter a cell only if $\textit{myTime}<\textit{monsterTime}$.
  \item If a boundary cell is reached, reconstruct and print the path
        (\texttt{YES}); else \texttt{NO}.
\end{enumerate}

\subsection*{Dry run}

\dryrun{\texttt{A} near the centre with one distant \texttt{M}.}

\begin{itemize}
  \item Monster BFS marks cells near \texttt{M} with small times, cells
        near \texttt{A} with large times.
  \item Your BFS heads to the nearest boundary, reaching it before the
        monster's time there -- escape with a printed path.
\end{itemize}

\subsection*{C++ implementation}

\begin{lstlisting}
#include <bits/stdc++.h>
using namespace std;

int main() {
    int n, m;
    cin >> n >> m;
    vector<string> g(n);
    for (auto& row : g) cin >> row;

    const int INF = 1e9;
    vector<vector<int>> mt(n, vector<int>(m, INF));   // monster arrival time
    queue<pair<int,int>> q;
    int sr = 0, sc = 0;
    for (int i = 0; i < n; i++)
        for (int j = 0; j < m; j++) {
            if (g[i][j] == 'M') { mt[i][j] = 0; q.push({i, j}); }
            if (g[i][j] == 'A') { sr = i; sc = j; }
        }
    int dr[] = {-1,1,0,0}, dc[] = {0,0,-1,1};
    char mvc[] = {'U','D','L','R'};

    while (!q.empty()) {                              // multi-source BFS
        auto [r, c] = q.front(); q.pop();
        for (int k = 0; k < 4; k++) {
            int nr = r + dr[k], nc = c + dc[k];
            if (nr < 0 || nr >= n || nc < 0 || nc >= m) continue;
            if (g[nr][nc] == '#' || mt[nr][nc] != INF) continue;
            mt[nr][nc] = mt[r][c] + 1;
            q.push({nr, nc});
        }
    }

    vector<vector<int>> myt(n, vector<int>(m, INF));
    vector<vector<char>> mv(n, vector<char>(m, 0));
    myt[sr][sc] = 0;
    q.push({sr, sc});
    int er = -1, ec = -1;
    while (!q.empty()) {
        auto [r, c] = q.front(); q.pop();
        if (r == 0 || r == n - 1 || c == 0 || c == m - 1) { er = r; ec = c; break; }
        for (int k = 0; k < 4; k++) {
            int nr = r + dr[k], nc = c + dc[k];
            if (nr < 0 || nr >= n || nc < 0 || nc >= m) continue;
            if (g[nr][nc] == '#' || myt[nr][nc] != INF) continue;
            if (myt[r][c] + 1 >= mt[nr][nc]) continue;    // monster reaches first
            myt[nr][nc] = myt[r][c] + 1;
            mv[nr][nc] = mvc[k];
            q.push({nr, nc});
        }
    }

    if (er == -1) { cout << "NO\n"; return 0; }
    string path;
    int r = er, c = ec;
    while (!(r == sr && c == sc)) {
        char ch = mv[r][c]; path += ch;
        if (ch == 'U') r++; else if (ch == 'D') r--;
        else if (ch == 'L') c++; else c--;
    }
    reverse(path.begin(), path.end());
    cout << "YES\n" << path.size() << "\n" << path << "\n";
    return 0;
}
\end{lstlisting}

\begin{complexitybox}
\textbf{Time Complexity.} $O(nm)$ -- two BFS passes over the grid.
\textbf{Space Complexity.} $O(nm)$ for the time and move grids.
\end{complexitybox}

\begin{takeawaybox}
Races against spreading hazards are multi-source BFS (the hazard field)
plus a guarded single-source BFS (your moves), entering a cell only when
you strictly beat the hazard. This two-field comparison is the template for
fire/flood-escape problems.
\end{takeawaybox}

